% Tests for latex-pro-calculus.sty, loaded on its own.
\documentclass{article}
\usepackage{latex-pro-calculus}
\usepackage{lptest}
\begin{document}
\lpsame{partialdiff first order}{\partialdiff[f]{x}}{\frac{\partial f}{\partial x}}
\lpsame{partialdiff mixed}{\partialdiff[f]{x,y}}{\frac{\partial^{2} f}{\partial x\partial y}}
\lpsame{partialdiff repeats collapse}{\partialdiff[f]{x,x,y}}
  {\frac{\partial^{3} f}{\partial x^{2}\partial y}}
\lpsame{partialdiff keeps written order}{\partialdiff[f]{x,y,x}}
  {\frac{\partial^{3} f}{\partial x\partial y\partial x}}
\lpsame{partialdiff operator}{\partialdiff[]{t}}{\frac{\partial}{\partial t}}
% \partialdiff without [function] uses the default function f
\lpsame{partialdiff default function}{\partialdiff{y}}{\frac{\partial f}{\partial y}}
\lpsame{partialdiff default same as [f]}{\partialdiff{x,x,y}}{\partialdiff[f]{x,x,y}}
\lpsame{partialdiff explicit function wins}{\partialdiff[g]{y}}{\frac{\partial g}{\partial y}}
\setpartialdifffunction{g}
\lpsame{setpartialdifffunction}{\partialdiff{y}}{\frac{\partial g}{\partial y}}
\lpsame{setpartialdifffunction higher order}{\partialdiff{x,y}}
  {\frac{\partial^{2} g}{\partial x\partial y}}
\lpsame{setpartialdifffunction explicit wins}{\partialdiff[h]{y}}{\frac{\partial h}{\partial y}}
\lpsame{setpartialdifffunction empty bracket still operator}{\partialdiff[]{y}}
  {\frac{\partial}{\partial y}}
\lpsame{setpartialdifffunction does not change hessian}{\hessian*{1}}
  {\mat{\dfrac{\partial^{2} f}{\partial x_1^{2}}}}
\setpartialdifffunction{u(x,y)}
\lpsame{setpartialdifffunction with arguments}{\partialdiff{x}}
  {\frac{\partial u(x,y)}{\partial x}}
{\setpartialdifffunction{h}}
\lpsame{setpartialdifffunction is local to a group}{\partialdiff{x}}
  {\frac{\partial u(x,y)}{\partial x}}
\setpartialdifffunction{}
\lpsame{setpartialdifffunction empty gives operator}{\partialdiff{t}}
  {\frac{\partial}{\partial t}}
\setpartialdifffunction{f}
\lpsame{diff upright d}{\diff[f]{x,x}}{\frac{\mathrm{d}^{2} f}{\mathrm{d} x^{2}}}
\lpsame{diff does not leak into partialdiff}{\diff{x}\partialdiff[]{x}}
  {\frac{\mathrm{d}}{\mathrm{d} x}\frac{\partial}{\partial x}}
\lpsame{diff keeps no default function}{\diff{x}}{\frac{\mathrm{d}}{\mathrm{d} x}}


\lpsame{hessian 2}{\hessian{2}}
  {\mat{\dfrac{\partial^{2} f(v)}{\partial x_1^{2}} & \dfrac{\partial^{2} f(v)}{\partial x_1\partial x_2} \\
        \noalign{\vskip 1ex}
        \dfrac{\partial^{2} f(v)}{\partial x_2\partial x_1} & \dfrac{\partial^{2} f(v)}{\partial x_2^{2}}}}
% regression: empty options keep their defaults
\lpsame{hessian empty first option}{\hessian[][g]{5}}{\hessian[x][g][v]{5}}
\lpsame{hessian empty first two options}{\hessian[][][\mathbf a]{2}}{\hessian[x][f][\mathbf a]{2}}
\lpsame{hessian all options empty}{\hessian[][][]{3}}{\hessian{3}}
\lpsame{hessian star drops the point}{\hessian*{2}}
  {\mat{\dfrac{\partial^{2} f}{\partial x_1^{2}} & \dfrac{\partial^{2} f}{\partial x_1\partial x_2} \\
        \noalign{\vskip 1ex}
        \dfrac{\partial^{2} f}{\partial x_2\partial x_1} & \dfrac{\partial^{2} f}{\partial x_2^{2}}}}
\lpsame{hessian variable list}{\hessian[x,y]{}}
  {\mat{\dfrac{\partial^{2} f(v)}{\partial x^{2}} & \dfrac{\partial^{2} f(v)}{\partial x\partial y} \\
        \noalign{\vskip 1ex}
        \dfrac{\partial^{2} f(v)}{\partial y\partial x} & \dfrac{\partial^{2} f(v)}{\partial y^{2}}}}
\lpsame{hessian list ignores size}{\hessian[x,y]{2}}{\hessian[x,y]{}}

\lpsame{gradient}{\gradient{2}}
  {\mat{\dfrac{\partial f(v)}{\partial x_1} & \dfrac{\partial f(v)}{\partial x_2}}}
\lpsame{gradient empty option}{\gradient[][h]{2}}{\gradient[x][h][v]{2}}
\lpsame{gradient list}{\gradient*[x,y]{}}
  {\mat{\dfrac{\partial f}{\partial x} & \dfrac{\partial f}{\partial y}}}
\lpsame{jacobian}{\jacobian*{2}{1}}
  {\mat{\dfrac{\partial f_1}{\partial x_1} \\ \noalign{\vskip 1ex} \dfrac{\partial f_2}{\partial x_1}}}
\lpsame{jacobian empty option}{\jacobian[][g]{2}{2}}{\jacobian[x][g][v]{2}{2}}
\lpsame{calculus follows bracket style}{\usematb\gradient*{1}}{\bmat{\dfrac{\partial f}{\partial x_1}}}

\begin{equation}\hessian[][g]{2}\label{eq:h}\end{equation}
\lpref{label after hessian}{eq:h}{1}
\end{document}
